\documentclass[a4paper,12pt]{article}
\usepackage{graphicx, setspace, array, xcolor, mathtools, contour, tikz, amsthm, setspace, amsmath,amsfonts,amssymb, subcaption, fancyhdr, lipsum,multicol, geometry, gensymb, hyperref}
\usepackage{colortbl}
\usepackage{bm}
\geometry{a4paper, margin=0.75in}
\contourlength{0.5pt}
\usetikzlibrary{patterns}
\usepackage[english]{babel}
\renewcommand{\qedsymbol}{$\blacksquare$}
\definecolor{darkgreen}{rgb}{0.0, 0.5, 0.0}
\pagestyle{fancy}
\fancyhf{}
\renewcommand{\headrulewidth}{0.4pt}
\fancyhead[L]{\rightmark}
\fancyhead[R]{\thepage}
\title{Modern Algebra}
\author{Assignment 14\\ \\ Yousef A. Abood\\ \\ ID: 900248250}
\date{December 2025}
\setlength{\parindent}{0pt}
\singlespacing
\parskip=1mm
\setcounter{section}{13}
\setcounter{subsection}{1}
\onehalfspacing
\begin{document}
\maketitle
\noindent\makebox[\linewidth]{\rule{15cm}{0.4pt}}
\section{Ideals}
\section*{Problem 1}
\begin{proof}
    Suppose we have a commutative ring with unity $R$, let $\langle a \rangle$ be the set $\{ra \mid r \in R\}.$
    We need to show that $\langle a \rangle$ is an ideal. First, we show that it is closed under subtraction. Consider the elements $ra, r'a \in \langle a \rangle$, where $r, r' \in R$, and observe that $ra-r'a=(r-r')a$ by multiplication distributivity. Since $R$ is closed under addition and $r, -r' \in R$, so $(r-r')=(r+(-r')) \in R.$ Thus, $(r+(-r'))a \in \langle a \rangle$. Hence, $\langle a \rangle$ is closed under subtraction.\\ 
    Next, we need to show that $ar, ra \in \langle a \rangle$ for any $r \in R.$ Consider an element $x \in \langle a \rangle$, by definition of $\langle a \rangle$, $x =r'a$ for some $r'$ in $R.$ Let $s$ be an element in $R$ and observe that $sx=s(r'a)=(sr')a$. Clearly, $sr' \in R$ so $sx=(sr')a \in \langle a \rangle.$ Hence, $ra \in \langle a \rangle$ for any $r \in R$. Since $R$ is commutative and $a \in R$, $ar=ra$. Thus, $ra, ar \in \langle a \rangle.$ Therefore, $\langle a \rangle$ is an ideal.     
\end{proof}
\section*{Problem 3}
%  Verify that the set I in Example 5 is an ideal and that if J is any
% ideal of R that contains a 1, a2, . . . , an, then I # J. (Hence, ka1,
% a2, . . . , anl is the smallest ideal of R that contains a1, a2, . . . , an .)
% EXAMPLE 5 Let R be a commutative ring with unity and let a1,
% a2, . . . , an belong to R. Then I 5 ka1, a2, . . . , anl 5 {r 1a1 1 r 2a2 1
% ? ? ? 1 r nan | ri [ R} is an ideal of R called the ideal generated by a1,
% a2, . . . , an. The verification that I is an ideal is left as an easy exercise
\begin{proof}
    Suppose we have a commutative ring with unity $R$, let $\langle a_1,a_2, \cdots, a_n \rangle$ be the set $$I=\{r_1a_1+ r_2a_2+ \cdots+ r_na_n \mid r_i \in R\}.$$
    We need to show that $I$ is an ideal. First, we show that it is closed under subtraction. Consider the elements $\sum_{i=1}^{n}r_ia_i, \sum_{i=1}^{n}r'_ia_i \in I$, where $r_i, r'_i \in R$, and observe that \begin{align*}\sum_{i=1}^{n}r_ia_i- \sum_{i=1}^{n}r'_ia_i&=\sum_{i=1}^{n}r_ia_i-r'_ia_i\\ &=\sum_{i=1}^{n}(r_i-r'_i)a_i \\ &= (r_1-r'_1)a_1+(r_2-r'_2)a_2+\cdots+(r_n-r'_n)a_n\end{align*} by multiplication distributivity. Since $R$ is closed under addition and $r_i, -r_i' \in R$, so $(r_i-r'_i)=(r_i+(-r'_i)) \in R.$ Thus, $\sum_{i=1}^{n}(r_i+(-r'_i))a_i \in I$. Hence, $I$ is closed under subtraction.\\ 
    Next, we need to show that $\sum_{i=1}^{n}a_ir_i, \sum_{i=1}^{n}r_ia_i \in I$ for any $r_i \in R.$ Consider an element $x \in I$, by definition of $I$, $x =\sum_{i=1}^{n}r'_ia_i$ for some $r'_i$ in $R.$ Let $s$ be an element in $R$ and observe that
    \begin{align*}
        sx=s\sum_{i=1}^{n}(r_i'a_i)=\sum_{i=1}^{n}s(r'_ia_i)=\sum_{i=1}^{n}(sr'_i)a_i.
    \end{align*}
    Clearly, $sr'_i \in R$ so $sx=\sum_{i=1}^{n}(sr'_i)a_i \in I.$ Hence, $\sum_{i=1}^{n}r_ia_i \in I$ for any $r_i \in R$. Since $R$ is commutative and $a \in R$, $\sum_{i=1}^{n}a_ir_i=\sum_{i=1}^{n}r_ia_i$. Thus, $\sum_{i=1}^{n}a_ir_i,\sum_{i=1}^{n}r_ia_i \in I.$ Therefore, $\langle a_1, a_2, \cdots , a_n \rangle$ is an ideal.     
\end{proof}
\section*{Problem 4}
\textbf{Consider the set $\bm{I = \{(n,n) \mid n \in \mathbb{Z}\}.}$}
\begin{proof}
 We need to show that it is a subring. First, pick $(n,n), (m,m) \in I$ where $n,m \in \mathbb{Z}.$ Observe that $(n,n)-(m,m)=(n-m,n-m)$. We know that $n-m \in \mathbb{Z}$ so $(n-m,n-m)$ is in $I.$ Next, we show that $I$ is closed under multiplication. See that $(n,n) \cdot (m,m)=(nm,nm)$, but $nm \in \mathbb{Z}$ so $(nm,nm) \in I.$ Hence, we showed that $I$ is a subring.
Now, we show that it is not an ideal. In other words, we show that $I$ violates one of ideal properties. Observe that when we multiply the element $(1,1) \in I$ by $(0,1) \in \mathbb{Z}\oplus\mathbb{Z}$ we get $(1,1)\cdot (0,1)=(1 \cdot 0, 1 \cdot 1)=(0,1) \notin I.$ Therefore, $I$ does not have the absorption property and it is not an ideal.   
\end{proof}
\section*{Problem 5}
\textbf{Consider the set $\bm{S = \{a+bi \mid a,b \in \mathbb{Z}, b \text{ \textit{is even}}\}.}$}
\begin{proof}
    First, we need to show that $S$ is a subring of $\mathbb{Z}[i].$ We begin by showing that $S$ is closed under subtraction. Let $a,b,c,d \in \mathbb{Z}$ where $b,d$ are even integers, observe that $(a+bi)-(c+di)=(a-c)+(b-d)i.$ Since $(a-c), (b-d) \in \mathbb{Z}$ and $b-d$ is even, $S$ is closed under subtraction. Now, we show that $S$ is closed under multiplication. Consider $a+bi,c+di \in S$, and observe that $(a+bi) \cdot (c+di)=ac+adi+cbi+bdi^{2}=(ac-bd)+(ad+cb)i$. Clearly, $(ac-bd), (ad+cb) \in \mathbb{Z}$ and $(ad+cb)$ is even. Hence, $S$ is closed under multiplication. Hence, $S$ is a subring.
    Next, we show that it is not an ideal as it violates the absorption property. That is, observe that $(1+2i) \cdot (3i)=3i+6i^2=3i+6(-1)=-6+3i$. Since $3$ is not even then $-6+3i \notin S$ and it is not an ideal.
\end{proof}
\section*{Problem 11}
\begin{itemize}
    \item [a)] We can write the set of elements in $\langle 2 \rangle + \langle 3 \rangle$ as $\{2a+3b \mid a,b \in \mathbb{Z}\}.$ Clearly, $2a+3b$ is a linear combination. Since the greatest common divisor of two elements divides all linear combinations of these two elements, then the set of all linear combinations is exactly the set generated by the $\operatorname{gcd}.$ Thus, $\langle 2 \rangle + \langle 3 \rangle = \langle 1 \rangle.$
    \item [b)] Similarly, $\langle 6 \rangle + \langle 8 \rangle = \langle \operatorname{gcd}(6,8) \rangle = \langle 2 \rangle.$
\end{itemize}
\phantomsection \label{sec:prob15}
\section*{Problem 15}
\begin{proof}
  Suppose we have a ring $R$ and an ideal $A$ of that ring, so $A \subseteq R$. Assume that $1$ belongs to $A.$ Let $r,a$ be elements in $R,A$ respectively. By definition of Ideal, we know that $ra \in A.$ Choose $a$ to be $1$ and $r$ to be any element in $R$.  By $r \cdot 1 =r$, every element in $R$ is in $A$, that is, $R \subseteq A.$ Therefore, since $A \subseteq R \wedge R \subseteq A$, $A=R.$  
\end{proof}
\section*{Problem 32}
\begin{proof}
  Consider the ring $R=\mathbb{Z}\oplus \mathbb{Z}$, let $A = \{(3x,y) \mid x,y \in \mathbb{Z}\}$ and $A \subseteq R.$ First, we need to show that $A$ is an ideal. We begin by showing that it is closed under subtraction, let $a,b,x,y \in \mathbb{Z}$ and consider the elements $(3a,b), (3x,y) \in A.$ Observe that $(3a,b)-(3x,y)=(3a-3x, b-y)=(3(a-x), b-y)$. Clearly, $b-y, a-x$ are integers, so $(3(a-x), b-y)\in A$. Thus, $A$ is closed under subtraction. Second, we show that it has the absorption property. Let $(r,s)$ be any element in $R$, where $r,s$ are integers. Observe that $(r,s) \cdot (3x,y)=(3xr, ys)$. Clearly, since $xr, ys$ are integers, $(3xr, ys) \in A$ and it satisfies the property. 
  \\ To show that $A$ is maximal assume that there exists an ideal $B$ of $R$ such that $A \subsetneq B \subseteq R.$ By our assumption, there is an element $\beta =(u,v) \in B\setminus A$ where $u,v \in \mathbb{Z}$, so $\beta \notin A.$
    Consider the element $(0,1)=(3\cdot 0, 1)$, so it is in $A$. Thus, $(0,1)$ is also in $B.$ Using ideal second property, observe that $(u,0)=(u,v)-v(0,1) \in B$, where $3 \nmid u.$ Using \textit{Bézout's Theorem}, we know that $\operatorname{gcd}(u,3)=1$ so we pick elements $m,n \in \mathbb{Z}$ such that $1=3m+un.$ Then $(1,0)=(3m+un,0)=(3m,0)+(un,0)=m(3,0)+n(u,0)$ and $(1,0) \in B.$ Since $(1,0), (0,1) \in B$, $(1,0)+(0,1)=(1,1)$ is also in $B$. We know that $(1,1)$ is the multiplicative identity of $R$. Hence, by our work in \hyperref[sec:prob15]{Problem 15}, $B=R.$ Therefore, $A$ is maximal.\\
    For the other part, suppose we replace $3x$ by $4x$ in our definition of $A.$ Trivially, we can find another ideal $S=\{(2x,y) \mid x,y \in \mathbb{Z}\}$ such that $A \subsetneq S.$ Thus, $A$ is not maximal. The reason why it was maximal in the other case is that $3$ is a prime number, so the greatest common divisor between it and any other element outside $A$ is $1$ which will help us to find that any maximal ideal $B$ such that $A \subsetneq B$ is the ring itself. Therefore, we generalize the statement:
    \begin{center}
        \textit{\textbf{An ideal $\bm{A=\{(nx,y) \mid x,y \in \mathbb{Z}\}}$ of the ring $\bm{\mathbb{Z}\oplus \mathbb{Z}}$ is maximal if $\bm{n}$ is prime.}}
    \end{center}
\end{proof}

\section*{Problem 37}
\begin{proof}
   Let $R$ be the ring $\mathbb{Z}\oplus \mathbb{Z}$ and $I=\{(a,0) \mid a \in \mathbb{Z}\}.$ We need to show that $I$ is a prime ideal. First, we show that $I$ is an ideal. We begin by showing closure under subtraction. See that $(x,0), (y,0) \in I$, and observe that $(x,0)-(y,0)=(x-y,0).$ Since $x-y$ is an integer, then $(x-y,0) \in I$. Thus, $I$ is closed under subtraction. Next, we show that it has the absorption property. Pick any element $(c,b) \in R,$ where $c,b \in \mathbb{Z}.$ Observe that $(c,b)\cdot (a,0)=(ca,0)$. Thus, $(ca,0) \in I$ since $ca \in \mathbb{Z}.$
\\ Now, we will show that $I$ is prime. Let $\beta=(u,v),\alpha =(s,t) \in R$, where $u,v,s,t \in \mathbb{Z}$. Observe $\beta \alpha = (us, vt)$. If $\beta \alpha \in I$, then $(us,vt)=(us, 0)$ and $vt=0.$ By definition of integer multiplication, either $v=0$ or $t=0$. Thus, we will have $\alpha=(s,0) \in I$ or $\beta = (u,0) \in I.$ Hence, $I$ is a prime ideal. 
\\To show it is not maximal, observe the set $S=\{(a,b) \mid a,b \in \mathbb{Z}, \text{b is even}\}.$ It is an ideal as it is clearly closed under subtraction as the subtraction of two integers is integer and the subtraction of two even numbers is even, and it has the absorption property since $b$ multiplied by any integer produce an even integer. Since $0$ is even, then $(a,0) \in S$ for any $a \in \mathbb{Z}.$ Clearly $S \ne I$ as it contains $(1,2)$ and $S \ne R$ as $R$ contains $(1,1).$ Therefore, $I$ is not maximal.
\end{proof}
\end{document}
